\section*{Appendix} 

% \section{Related Work}
% \label{app:related-work}

% \paragraph{Test-time scaling and candidate selection.}
% Test-time scaling has emerged as an effective paradigm for improving LLM performance without additional training, including longer reasoning traces ~\citep{wei_chain--thought_nodate}, self-refinement~\citep{shinn_reflexion_nodate}, search, and best-of-$N$ sampling. These methods~\citep{agrawal_gepa_2025, novikov_alphaevolve_2025, brown2024largelanguagemonkeysscaling, snell_scaling_2024} improve generation by allocating more inference compute to either a single trajectory or a larger pool of candidate solutions. However, they often assume access to a reliable evaluator, such as a deterministic checker, majority vote, or coarse LM judge. In contrast, our work studies verification itself as the object of scaling. We show that candidate pools often contain substantial oracle headroom, and that downstream performance is bottlenecked by verification rather than by generation alone.

% \paragraph{LLM-as-a-judge and automatic evaluation.}
% LLM-as-a-judge methods provide a scalable alternative to human evaluation by prompting LLMs to assess generated outputs. Prior work such as MT-Bench and Chatbot Arena~\citep{zheng_judging_nodate} established the effectiveness of LLM judges for open-ended response evaluation, while G-Eval~\citep{liu_g-eval_2023} introduced chain-of-thought reasoning and probability-weighted scoring to improve alignment with human preferences. Weaver~\citep{saad2025shrinking} further shows that ensembling weak reward models can reduce the generation-verification gap. Our work builds on this line of work but differs in setting, objective, and scaling characterization. Rather than evaluating isolated natural-language responses, LLM-as-a-Verifier verifies long-horizon agent trajectories involving tool use, code execution, robotics, and medical decision-making. Moreover, instead of collapsing the model's scoring distribution into a single discrete rating, we use a probabilistic approach and study how verification quality scales with score granularity, repeated evaluations, and criteria decomposition. 

% \paragraph{Reward models, process supervision, and verifiers.}
% Learned reward models and verifiers~\citep{cobbe_training_2021, zhang2025generativeverifiersrewardmodeling} have been widely used to improve reasoning and decision-making. Outcome reward models score final solutions, while process reward models provide intermediate supervision over reasoning steps. In robotics, learned reward models and action verifiers~\citep{kwok_robomonkey_2025, kwok2026scalingverificationeffectivescaling, liu2026worldactionverifierselfimproving} have been used to score candidate actions or trajectories. However, these approaches typically require domain-specific training data and may not generalize across tasks with different modalities, horizons, or success criteria. LLM-as-a-Verifier instead uses a frozen language or vision-language model as a general-purpose trajectory reward model. The same approach can be applied across coding, robotics, and medical benchmarks without additional fine-tuning.

% \paragraph{Efficient pairwise ranking.}
% Selecting the best trajectory from a candidate pool can be done through pointwise scoring, majority voting, or exhaustive pairwise comparison. Full round-robin comparison provides strong ranking accuracy but requires $\mathcal{O}(N^2)$ verifier calls, which becomes expensive as the number of candidates grows. Recent pairwise verifier methods~\citep{singh_v_1_2026} reduce this cost, but often rely on coarse discrete judgments or fixed tournament schedules. Our Probabilistic Pivot Tournament uses continuous preference probabilities derived from verifier scores and concentrates comparisons on likely top candidates. A ring pass first mitigates positional bias, after which high-scoring candidates are selected as pivots for additional comparisons. This reduces the verification budget to $\mathcal{O}(Nk)$ while preserving much of the accuracy of full round-robin ranking.


% \vspace{-7pt}
% \paragraph{Test-Time Scaling.}


% Scaling test-time compute has emerged as a critical paradigm for enhancing LLM performance without additional training. We categorize these advancements into two primary dimensions: vertical and horizontal scaling. \textit{Vertical test-time scaling} enhances performance by increasing the depth or "thinking time" being allocated. Pioneered by Chain-of-Thought (CoT) prompting \citep{wei_chain--thought_nodate}, this paradigm has evolved to include sophisticated context management for long-term memory \citep{zhang_agentic_2026} and iterative self-correction loops such as Reflexion \citep{shinn_reflexion_nodate}. In contrast, \textit{horizontal test-time scaling} focuses on breadth, generating a diverse population of candidate responses and selecting the optimal output based on external or internal heuristics~\citep{agrawal_gepa_2025, novikov_alphaevolve_2025, brown2024largelanguagemonkeysscaling, snell_scaling_2024}. However, these methods often treat evaluation as fixed, either relying on deterministic task checkers, majority voting, or coarse model-based judges. In contrast, our work studies verification itself as the object of scaling. We show that when candidate solutions already contain substantial oracle headroom, downstream performance is bottlenecked by the ability to reliably select correct trajectories rather than by generation alone.

% \vspace{-7pt}
% \paragraph{LLM-as-a-Judge.}

% LLM-as-a-judge methods~\citep{zheng_judging_nodate} use language models as scalable evaluators for open-ended outputs, and probability-weighted variants such as G-Eval~\citep{liu_g-eval_2023} show that token probabilities can improve alignment with human judgments. Our work builds on this insight but differs in both setting and objective: rather than evaluating isolated natural-language responses, we verify long-horizon agent trajectories, where small differences in terminal state, tool use, and recovery behavior can determine success. Moreover, we explicitly characterize how verifier quality scales with scoring-token granularity, repeated evaluations, and criteria decomposition. (you should add MT-Bench and Chatbot Arena)


% \paragraph{Verification and Reward Models.}
% The concept of action verification—where multiple candidate actions are proposed and subsequently scored by a discriminator—has been proposed and applied in the field of robotics. Recent work in this direction include RoboMonkey \citep{kwok_robomonkey_2025}, which utilizes a trained verifier model to score actions, and TopReward \citep{chen_topreward_2026}, which leverages token probabilities as "hidden" zero-shot rewards to create lightweight quality estimators. Within the LLM agent domain, methods for action verification include "LLM-as-a-judge" pointwise scoring \citep{yuan_self-rewarding_2025}, and pairwise frameworks such as $V_1$ \citep{singh_v_1_2026} . However, existing verifiers often struggle in long-context scnearios. Our work extends these concepts by introducing weighted LLM logits as a reward model for agent action verification.

% % \paragraph{LLM-based Agents.}
% % Large Language Models (LLMs), pre-trained on internet-scale datasets, have demonstrated remarkable capabilities in zero-shot generalization, creative synthesis, and linguistic nuance. However, a significant performance gap remains in \textit{agentic} tasks: scenarios requiring long-horizon planning, precise tool manipulation, and autonomous error recovery within dynamic environments. These challenges are codified in increasingly difficult benchmarks such as SWE-bench \citep{jimenez_swe-bench_2024}, which focuses on real-world software engineering; Tau-bench \citep{yao_-bench_2024}, which evaluates tool-use consistency; and Terminal-bench \citep{merrill_terminal-bench_2026}, which tests proficiency in shell-based environments. Current research to improve agentic performance typically follows two trajectories: the development of architectural harnesses like ReAct \citep{yao_react_2023}, which interleaves reasoning with environmental interaction, and post-training alignment via reinforcement learning, such as Group Relative Policy Optimization (GRPO) \citep{shao_deepseekmath_2024}.